% GENERATED FILE. Edit canonical lessons, then run npm run generate:books.
% canonical-source-hash: fnv1a64:aa48e3b818abac11
% canonical-lessons: IT-C386-violare, IT-C386-violazione, IT-C386-infrazione, IT-C386-trasgredire, IT-C386-eccezione, IT-R386-ripasso-379-386

\chapter{Infrangere le regole}
\label{ch:it-infrangere-le-regole}
\hlchaptermodality{386}

\begin{chapteropening}
\textbf{By the end of this chapter:} \emph{I can say that someone broke a rule: a violation, an offence, crossing the line and an exception to the rule.}

Complete the last lesson of chapter 386: I can say that someone broke a rule: a violation, an offence, crossing the line and an exception to the rule.
\end{chapteropening}

\section[violare]{violare — to violate, to break (a rule)}

\label{lesson:IT-C386-violare}



\begin{quote}
Before the new one: say the Italian for to authorise, then the Italian for the authorisation.
\end{quote}

\subsection*{You'll want to know: violare}
\textbf{violare} — \textquotedblleft{}to violate, to break (a rule)\textquotedblright{}.

\begin{quote}
Chi viola le regole paga una multa.
\end{quote}

\begin{quote}
\emph{Anyone who breaks the rules pays a fine.}
\end{quote}

Stress: \emph{vìolo}.

\begin{cousinweb}
\emph{violare} is a learned borrowing from Latin \emph{violāre} \textquotedblleft{}to treat with violence, to break\textquotedblright{} (Lewis \& Short), related to \emph{vīs} \textquotedblleft{}force\textquotedblright{}.
\end{cousinweb}

\begin{grammarlens}[title={What you've built}]
One more word for telling, explaining and describing at B1.
\end{grammarlens}

\subsection*{Your turn}
\begin{itemize}
\raggedright
  \item \emph{Say it:} \emph{violare}
  \item \emph{Say it:} \emph{violare}, once more
  \item \emph{Say it:} a sentence of your own with \emph{violare}
  \item \emph{Recall:} say the Italian for to authorise, then the Italian for the authorisation, then say \emph{violare} again
\end{itemize}

\begin{tcolorbox}[breakable,colback=teal!4,colframe=teal!35!black,title={Before you move on}]
What does violare mean? (\textquotedblleft{}To violate, to break (a rule)\textquotedblright{}.) Say it once more.
\end{tcolorbox}
\section[la violazione]{la violazione — the violation, the breach}

\label{lesson:IT-C386-violazione}



\begin{quote}
Before the new one: say the Italian for the authorisation, then the Italian for to violate.
\end{quote}

\subsection*{You'll want to know: la violazione}
\textbf{la violazione} — \textquotedblleft{}the violation, the breach\textquotedblright{}.

\begin{quote}
Questa è una violazione della privacy.
\end{quote}

\begin{quote}
\emph{This is a breach of privacy.}
\end{quote}

\emph{In violazione di} — in breach of.

\begin{cousinweb}
\emph{violazione} is a learned borrowing from Latin \emph{violātiō} \textquotedblleft{}an injury, a profaning\textquotedblright{} (Lewis \& Short), from \emph{violāre} \textquotedblleft{}to treat with violence\textquotedblright{}.
\end{cousinweb}

\begin{grammarlens}[title={What you've built}]
One more word for telling, explaining and describing at B1.
\end{grammarlens}

\subsection*{Your turn}
\begin{itemize}
\raggedright
  \item \emph{Say it:} \emph{la violazione}
  \item \emph{Say it:} \emph{la violazione}, once more
  \item \emph{Say it:} a sentence of your own with \emph{la violazione}
  \item \emph{Recall:} say the Italian for the authorisation, then the Italian for to violate, then say \emph{la violazione} again
\end{itemize}

\begin{tcolorbox}[breakable,colback=teal!4,colframe=teal!35!black,title={Before you move on}]
What does la violazione mean? (\textquotedblleft{}The violation, the breach\textquotedblright{}.) Say it once more.
\end{tcolorbox}
\section[l'infrazione]{l'infrazione — the offence, the infringement (of a rule)}

\label{lesson:IT-C386-infrazione}



\begin{quote}
Before the new one: say the Italian for to violate, then the Italian for the violation.
\end{quote}

\subsection*{You'll want to know: l'infrazione}
\textbf{l'infrazione} — \textquotedblleft{}the offence, the infringement (of a rule)\textquotedblright{}.

\begin{quote}
Passare con il rosso è un'infrazione grave.
\end{quote}

\begin{quote}
\emph{Going through a red light is a serious offence.}
\end{quote}

A breach of a rule rather than a crime: traffic and sport rules.

\begin{cousinweb}
\emph{infrazione} is a learned borrowing from Latin \emph{īnfrāctiō} \textquotedblleft{}a breaking\textquotedblright{} (Lewis \& Short), from \emph{īnfringere} \textquotedblleft{}to break in\textquotedblright{}, built on \emph{frangere} \textquotedblleft{}to break\textquotedblright{}.
\end{cousinweb}

\begin{grammarlens}[title={What you've built}]
One more word for telling, explaining and describing at B1.
\end{grammarlens}

\subsection*{Your turn}
\begin{itemize}
\raggedright
  \item \emph{Say it:} \emph{l'infrazione}
  \item \emph{Say it:} \emph{l'infrazione}, once more
  \item \emph{Say it:} a sentence of your own with \emph{l'infrazione}
  \item \emph{Recall:} say the Italian for to violate, then the Italian for the violation, then say \emph{l'infrazione} again
\end{itemize}

\begin{tcolorbox}[breakable,colback=teal!4,colframe=teal!35!black,title={Before you move on}]
What does l'infrazione mean? (\textquotedblleft{}The offence, the infringement (of a rule)\textquotedblright{}.) Say it once more.
\end{tcolorbox}
\section[trasgredire]{trasgredire — to break (a rule), to transgress}

\label{lesson:IT-C386-trasgredire}



\begin{quote}
Before the new one: say the Italian for the violation, then the Italian for the offence.
\end{quote}

\subsection*{You'll want to know: trasgredire}
\textbf{trasgredire} — \textquotedblleft{}to break (a rule), to transgress\textquotedblright{}.

\begin{quote}
Da ragazzo mi piaceva trasgredire le regole.
\end{quote}

\begin{quote}
\emph{As a boy I liked breaking the rules.}
\end{quote}

Conjugated with \emph{-isc-}: \emph{trasgredisco}. \emph{Trasgressivo} — provocative, rule-breaking.

\begin{cousinweb}
\emph{trasgredire} is a learned borrowing from Latin \emph{trānsgredī} \textquotedblleft{}to step across\textquotedblright{} (Lewis \& Short), built from \emph{trāns} \textquotedblleft{}across\textquotedblright{} and \emph{gradī} \textquotedblleft{}to step\textquotedblright{}: you step over the line.
\end{cousinweb}

\begin{grammarlens}[title={What you've built}]
One more word for telling, explaining and describing at B1.
\end{grammarlens}

\subsection*{Your turn}
\begin{itemize}
\raggedright
  \item \emph{Say it:} \emph{trasgredire}
  \item \emph{Say it:} \emph{trasgredire}, once more
  \item \emph{Say it:} a sentence of your own with \emph{trasgredire}
  \item \emph{Recall:} say the Italian for the violation, then the Italian for the offence, then say \emph{trasgredire} again
\end{itemize}

\begin{tcolorbox}[breakable,colback=teal!4,colframe=teal!35!black,title={Before you move on}]
What does trasgredire mean? (\textquotedblleft{}To break (a rule), to transgress\textquotedblright{}.) Say it once more.
\end{tcolorbox}
\section[l'eccezione]{l'eccezione — the exception}

\label{lesson:IT-C386-eccezione}



\begin{quote}
Before the new one: say the Italian for the offence, then the Italian for to break.
\end{quote}

\subsection*{You'll want to know: l'eccezione}
\textbf{l'eccezione} — \textquotedblleft{}the exception\textquotedblright{}.

\begin{quote}
Di solito non mangio dolci, ma oggi faccio un'eccezione.
\end{quote}

\begin{quote}
\emph{I do not usually eat sweets, but today I am making an exception.}
\end{quote}

\emph{Ad eccezione di} — except for. \emph{Senza eccezioni} — with no exceptions.

\begin{cousinweb}
\emph{eccezione} is a learned borrowing from Latin \emph{exceptiō} \textquotedblleft{}a taking out, an exception\textquotedblright{} (EtymDB), from \emph{excipere} \textquotedblleft{}to take out\textquotedblright{}, built on \emph{capere} \textquotedblleft{}to take\textquotedblright{}.
\end{cousinweb}

\begin{grammarlens}[title={What you've built}]
One more word for telling, explaining and describing at B1.
\end{grammarlens}

\subsection*{Your turn}
\begin{itemize}
\raggedright
  \item \emph{Say it:} \emph{l'eccezione}
  \item \emph{Say it:} \emph{l'eccezione}, once more
  \item \emph{Say it:} a sentence of your own with \emph{l'eccezione}
  \item \emph{Recall:} say the Italian for the offence, then the Italian for to break, then say \emph{l'eccezione} again
\end{itemize}

\begin{tcolorbox}[breakable,colback=teal!4,colframe=teal!35!black,title={Before you move on}]
What does l'eccezione mean? (\textquotedblleft{}The exception\textquotedblright{}.) Say it once more.
\end{tcolorbox}
\section[(dialogue)]{Ripasso dei capitoli 379–386}

\label{lesson:IT-R386-ripasso-379-386}



\begin{quote}
Say the Italian for the upper secondary school and to attend.
\end{quote}

\subsection*{Your turn}
Say these aloud:

\begin{itemize}
\raggedright
  \item the upper secondary school, to attend, the absence, the school report, the head teacher
  \item the faculty, the first-year student, to enrol, the enrolment, the semester
  \item the teacher, the thesis supervisor, the scholarship, the doctorate, the exam session
  \item to revise, the oral test, to memorise, the notes, the textbook
  \item to fail, the failing, to pass, the assessment, the score
  \item the rule, the regulations, compulsory, optional, the obligation
  \item forbidden, to ban, to tolerate, to authorise, the authorisation
  \item to violate, the violation, the offence, to break, the exception
\end{itemize}

\begin{tcolorbox}[breakable,colback=teal!4,colframe=teal!35!black,title={Before you move on}]
The upper secondary school? (\textbf{il liceo}.) To attend? (\textbf{frequentare}.) The absence? (\textbf{l'assenza}.) The school report? (\textbf{la pagella}.) The head teacher? (\textbf{il preside}.) The faculty? (\textbf{la facoltà}.) The first-year student? (\textbf{la matricola}.) To enrol? (\textbf{iscriversi}.) The enrolment? (\textbf{l'iscrizione}.) The semester? (\textbf{il semestre}.) The teacher? (\textbf{il docente}.) The thesis supervisor? (\textbf{il relatore}.) The scholarship? (\textbf{la borsa di studio}.) The doctorate? (\textbf{il dottorato}.) The exam session? (\textbf{l'appello}.) To revise? (\textbf{ripassare}.) The oral test? (\textbf{l'interrogazione}.) To memorise? (\textbf{memorizzare}.) The notes? (\textbf{gli appunti}.) The textbook? (\textbf{il manuale}.) To fail? (\textbf{bocciare}.) The failing? (\textbf{la bocciatura}.) To pass? (\textbf{promuovere}.) The assessment? (\textbf{la valutazione}.) The score? (\textbf{il punteggio}.) The rule? (\textbf{la norma}.) The regulations? (\textbf{il regolamento}.) Compulsory? (\textbf{obbligatorio}.) Optional? (\textbf{facoltativo}.) The obligation? (\textbf{l'obbligo}.) Forbidden? (\textbf{vietato}.) To ban? (\textbf{proibire}.) To tolerate? (\textbf{tollerare}.) To authorise? (\textbf{autorizzare}.) The authorisation? (\textbf{l'autorizzazione}.) To violate? (\textbf{violare}.) The violation? (\textbf{la violazione}.) The offence? (\textbf{l'infrazione}.) To break? (\textbf{trasgredire}.) The exception? (\textbf{l'eccezione}.)
\end{tcolorbox}
